\chapter{数据库：JPA 实体、复合主键、Repository}

这一关把站点和可用量存进数据库。这是 Spring Boot 里内容最多、也最容易卡住的一块，慢慢来。

\partinfo{%
  \code{station/persistence/StationEntity.java}：4.1\\
  \code{station/persistence/AvailabilityId.java}、\code{AvailabilityEntity.java}：4.2\\
  \code{station/persistence/StationRepository.java}：4.3\\
  \code{station/persistence/AvailabilityRepository.java}：4.4、4.5%
}{\code{./mvnw test -Dtest='Part4*'}}

\section{知识点}

\subsection{几个名字的关系}

\begin{itemize}
  \item \textbf{JPA}（Jakarta Persistence API）：Java 官方定义的一套“把对象存进关系数据库”的标准，主要是一组注解和接口，包名是 \code{jakarta.persistence}。
  \item \textbf{Hibernate}：JPA 最常用的实现。真正生成 SQL、执行查询的是它。
  \item \textbf{Spring Data JPA}：在 JPA 上面再包一层，让你只写接口、不写实现就能完成大部分查询。
\end{itemize}

\subsection{表结构由 Flyway 管理}

先打开 \code{src/main/resources/db/migration/V1__create_station_tables.sql} 看一下两张表：

\begin{sql}
CREATE TABLE station (
    number      INT          NOT NULL,
    name        VARCHAR(255) NOT NULL,
    address     VARCHAR(255),
    lat         DOUBLE       NOT NULL,
    lng         DOUBLE       NOT NULL,
    banking     BOOLEAN      NOT NULL,
    bike_stands INT          NOT NULL,
    status      VARCHAR(32),
    PRIMARY KEY (number)
);

CREATE TABLE availability (
    number                INT          NOT NULL,
    last_update           TIMESTAMP(6) NOT NULL,
    available_bikes       INT          NOT NULL,
    available_bike_stands INT          NOT NULL,
    status                VARCHAR(32),
    PRIMARY KEY (number, last_update),
    CONSTRAINT fk_availability_station FOREIGN KEY (number) REFERENCES station (number)
);
\end{sql}

应用启动时，Flyway 会按版本号依次执行 \code{db/migration} 下还没执行过的脚本。表结构因此有了版本记录，和代码一起提交到 git。

\code{application.yml} 里设置了 \code{spring.jpa.hibernate.ddl-auto: validate}，意思是：Hibernate 不建表、不改表，只\textbf{检查}你的实体类和数据库里的表是否对得上，对不上就启动失败。所以你的注解写错了，测试会直接告诉你哪一列有问题。

\subsection{实体（Entity）：一个类对应一张表}

\begin{java}
@Entity                               // 这是一个实体
@Table(name = "book")                 // 对应 book 表（不写的话默认用类名）
public class BookEntity {

    @Id                               // 主键
    private long isbn;

    @Column(nullable = false)         // 这一列不允许为空
    private String title;

    @Column(name = "page_count")      // 指定列名
    private int pageCount;

    private String author;            // 不加注解也会映射，列名默认按字段名

    protected BookEntity() {}         // JPA 要求有一个无参构造函数
    ...
}
\end{java}

\begin{itemize}
  \item 注解都在 \code{jakarta.persistence} 包里。
  \item Spring Boot 默认会把 camelCase 字段名转成 snake\_case 列名（\code{pageCount} → \code{page_count}），所以 \code{@Column(name = ...)} 很多时候可以省略。这一关要求你显式写出来，练习一下。
  \item JPA 通过反射创建对象，所以需要一个无参构造函数。写成 \code{protected} 是为了不让业务代码随便用它创建出不完整的对象。
\end{itemize}

\subsection{复合主键}

\code{availability} 表的主键有两列。JPA 的做法是：把两列包装成一个“可嵌入的值对象”，再把它作为实体的主键：

\begin{java}
@Embeddable                                          // 可以嵌入到实体里
public record LoanId(long isbn, @Column(name = "loan_date") LocalDate loanDate)
        implements Serializable {}                   // JPA 规范要求主键类可序列化

@Entity
public class LoanEntity {
    @EmbeddedId                                      // 这个字段是复合主键
    private LoanId id;
    ...
}
\end{java}

主键类还必须正确实现 \code{equals} 和 \code{hashCode}，用 record 的话这两个方法是自动生成的。

\begin{note}
Hibernate 7 其实很宽松：只要字段上有 \code{@EmbeddedId}，主键类上不加 \code{@Embeddable} 也能运行，测试也发现不了。但 JPA 规范要求加，换一个 JPA 实现可能就不行了，所以还是要写上。
\end{note}

\subsection{Repository：只写接口}

\begin{java}
public interface BookRepository extends JpaRepository<BookEntity, Long> {
    //                                               实体类型    主键类型
}
\end{java}

继承 \code{JpaRepository} 之后，不用写任何代码就有了 \code{save}、\code{saveAll}、\code{findById}、\code{findAll}、\code{count}、\code{existsById}、\code{deleteById} 等方法。Spring 在启动时自动生成实现类。

\textbf{派生查询（derived query）}：按规则给方法起名，Spring 会根据方法名自动生成查询：

\begin{java}
List<BookEntity> findByAuthor(String author);
// -> select ... where author = ?

List<BookEntity> findByAuthorOrderByTitleAsc(String author);
// -> select ... where author = ? order by title asc

Optional<BookEntity> findFirstByAuthorOrderByPublishedDesc(String author);
// -> 某个作者最新出版的一本书
\end{java}

方法名里的属性名是\textbf{实体的字段名}，不是表的列名。属性在嵌入对象里面时，把路径连起来写：主键 \code{id} 里面的 \code{isbn} 写成 \code{IdIsbn}，比如 \code{findByIdIsbn(long isbn)}。

\textbf{接口里的 default 方法}：Java 8 起，接口里可以写有方法体的 \code{default} 方法。这一关里，测试调用的是 \code{findOpenStations()}、\code{findLatest(...)} 这些名字好读的方法，你在里面调用自己声明的派生查询。

\subsection{自己写查询：@Query 和 JPQL}

方法名表达不了的查询，用 \code{@Query} 写 \textbf{JPQL}。JPQL 长得像 SQL，但写的是\textbf{实体名和字段名}，不是表名和列名：

\begin{java}
@Query("select b from BookEntity b where b.price > :minPrice")
List<BookEntity> findExpensive(@Param("minPrice") double minPrice);
\end{java}

一个常见需求是“每组取最新的一条”，比如每个作者最新出版的书。可以用子查询：

\begin{java}
@Query("""
        select b from BookEntity b
        where b.published = (
            select max(b2.published) from BookEntity b2
            where b2.author = b.author)
        """)
List<BookEntity> findNewestBookOfEachAuthor();
\end{java}

（\code{"""} 是 Java 的文本块，可以写多行字符串。）

\subsection{测试是怎么验证的}

\code{@DataJpaTest} 只启动数据库相关的部分：内存数据库 H2、Flyway、JPA 和 Repository。每个测试方法结束后会\textbf{自动回滚}，所以测试之间的数据互不影响。

测试里的 \code{flushAndClear()}：\code{flush} 把还没写进数据库的修改真正执行成 SQL，\code{clear} 清空 JPA 的一级缓存。这样接下来的查询一定会真的去查数据库，而不是直接从内存里拿刚才 save 的那个对象。

\section{题目}

\begin{task}{4.1\quad StationEntity}
给 \code{StationEntity} 加上注解，对应 \code{station} 表：
\begin{itemize}
  \item 类是一个实体，对应 \code{station} 表。
  \item \code{number} 是主键。
  \item \code{name} 列不允许为空。
  \item \code{bikeStands} 显式指定列名 \code{bike_stands}。
\end{itemize}

\begin{example}
\code{save} 一个站点，\code{flushAndClear} 之后用 \code{findById(42)} 读出来，所有字段都和存进去时一样。
\end{example}
\end{task}

\begin{task}{4.2\quad AvailabilityEntity 和复合主键}
\begin{itemize}
  \item \code{AvailabilityId} 是可嵌入的值对象；\code{lastUpdate} 对应 \code{last_update} 列。
  \item \code{AvailabilityEntity} 是实体，对应 \code{availability} 表；\code{id} 字段是复合主键；\code{availableBikes}、\code{availableBikeStands} 分别对应 \code{available_bikes}、\code{available_bike_stands} 列。
\end{itemize}

\begin{example}
\begin{plaincode}
save (42, 08:00) bikes=3
save (42, 08:05) bikes=5
findById(new AvailabilityId(42, 08:05))  ->  bikes=5
count()  ->  2

再 save (42, 08:00) bikes=4      // 主键相同
count()  ->  2                   // 不会多出一行，而是更新原来那行
\end{plaincode}
\end{example}
\end{task}

\begin{task}{4.3\quad 查询开放的站点}
在 \code{StationRepository} 里：

\textbf{a)} 声明一个派生查询方法：按 \code{status} 查找，结果按 \code{name} 升序。

\textbf{b)} \code{findOpenStations()} 调用它，\code{status} 传 \code{"OPEN"}。

\begin{example}
\begin{plaincode}
站点：CHARLEMONT PLACE (OPEN), BLESSINGTON STREET (CLOSED), ANNESLEY BRIDGE (OPEN)
findOpenStations()  ->  [ANNESLEY BRIDGE, CHARLEMONT PLACE]
\end{plaincode}
\end{example}
\end{task}

\begin{task}{4.4\quad 某个站点最新的一条记录}
在 \code{AvailabilityRepository} 里：

\textbf{a)} 声明一个派生查询方法：按主键里的 \code{number} 查找，按主键里的 \code{lastUpdate} 倒序，只取第一条，返回 \code{Optional}。

\textbf{b)} \code{findLatest(number)} 调用它。

\begin{example}
\begin{plaincode}
站点 42 有 08:00、08:05、08:10 三条记录
findLatest(42)  ->  08:10 那一条
findLatest(5)   ->  Optional.empty()（站点 5 没有记录）
\end{plaincode}
\end{example}
\end{task}

\begin{task}{4.5\quad 每个站点最新的一条记录}
把 \code{findLatestForAllStations()} 改成用 \code{@Query} 实现：\textbf{删掉 \code{default} 关键字和方法体}，在方法上加 \code{@Query}，写一条 JPQL。

\begin{example}
\begin{plaincode}
42: 08:00, 08:10
5:  08:00, 08:05
findLatestForAllStations()  ->  [(42, 08:10), (5, 08:05)]（顺序不限）
\end{plaincode}
\end{example}

Part 6 会用这个查询：一次查出所有站点的最新数据，而不是每个站点查一次。
\end{task}

\section{常见报错}

\begin{center}
\begin{tabularx}{\textwidth}{>{\raggedright\arraybackslash}X >{\raggedright\arraybackslash}X}
\toprule
报错信息里有 & 原因 \\
\midrule
\code{Not a managed type: class ...StationEntity} & 类上没加 \code{@Entity} \\
\code{Schema validation: missing column [xxx] in table [yyy]} & 字段对应的列名和表里的不一样 \\
\code{Schema validation: missing table [xxx]} & \code{@Table(name = ...)} 写错了 \\
\code{Entity '...' has no identifier} & 实体里没有标 \code{@Id} 或 \code{@EmbeddedId} 的字段 \\
\code{No property 'xxx' found for type 'yyy'} & 派生查询的方法名里，属性名拼错了 \\
\code{Could not resolve root entity 'xxx'} & \code{@Query} 的 JPQL 里写的是表名，应该写实体类名 \\
\bottomrule
\end{tabularx}
\end{center}

报错信息通常很长，而且同一个测试类里只要有一个测试启动失败，其他测试都会报 \code{ApplicationContext failure threshold exceeded}。\textbf{只看第一个失败的测试}，找 \code{Caused by:} 开头的最后一行，那里一般是真正的原因。

\begin{think}
\begin{enumerate}
  \item 生产环境为什么用 Flyway 管理表结构，而不是让 Hibernate 自动建表（\code{ddl-auto: update}）？
  \item 什么是 N+1 查询问题？如果 Part 6 里对每个站点各调用一次 \code{findLatest}，会发出多少条 SQL？
  \item \code{save()} 一个主键已经存在的实体时，JPA 是怎么决定执行 INSERT 还是 UPDATE 的？
  \item 为什么实体类不用 record？（提示：JPA 需要无参构造函数，而且要能修改字段）
\end{enumerate}
\end{think}
